% Preamble template for Tufte-style applied econometrics lecture notes.
% Do not compile this file directly; compile the main notes file instead.
\documentclass{tufte-handout}

\usepackage[T1]{fontenc}
\usepackage[utf8]{inputenc}
\usepackage{ntheorem}
\usepackage{graphicx}
\usepackage{amsmath}
\usepackage{amssymb}
\usepackage{mathtools}
\usepackage{hyperref}
\usepackage{epigraph}
\usepackage{booktabs}
\usepackage{array}
\theoremstyle{break}

\hypersetup{
  colorlinks,
  urlcolor = blue,
  pdfauthor={Swapnil Singh},
  pdfkeywords={econometrics, causal inference, applied econometrics},
  pdftitle={Lecture Notes for Applied Econometrics},
  pdfpagemode=UseNone
}

\newtheorem{ruleN}{Rule}
\newtheorem{thmN}{Theorem}
\newtheorem{assN}{Assumption}
\newtheorem{defN}{Definition}
\newtheorem{exmp}{Example}
\newtheorem{cmt}{Comment}
\newtheorem{proof}{Proof}

\newcommand{\continuation}{??}
\newtheorem*{excont}{Example \continuation}
\newenvironment{continueexample}[1]
 {\renewcommand{\continuation}{\ref{#1}}\excont[continued]}
 {\endexcont}

\newcommand{\bY}{\mathbf{Y}}
\newcommand{\bX}{\mathbf{X}}
\newcommand{\bD}{\mathbf{D}}
\newcommand{\E}{\mathbb{E}}
\newcommand{\Prob}{\mathbb{P}}
\newcommand{\Var}{\mathrm{Var}}
\newcommand{\Cov}{\mathrm{Cov}}
\newcommand{\se}{\mathrm{SE}}
\newcommand{\CI}{\mathrm{CI}}
\newcommand{\ATE}{\mathrm{ATE}}
\newcommand{\ATT}{\mathrm{ATT}}
\newcommand{\TOT}{\mathrm{TOT}}
\newcommand{\LATE}{\mathrm{LATE}}
\newcommand{\ITT}{\mathrm{ITT}}
\newcommand{\RD}{\mathrm{RD}}
\newcommand{\DiD}{\mathrm{DiD}}
\newcommand{\MMSE}{\mathrm{MMSE}}
\newcommand{\OLS}{\mathrm{OLS}}
\newcommand{\Normal}{\mathcal{N}}
\newcommand{\plim}{\operatorname*{plim}}
\DeclareMathOperator*{\argmin}{arg\,min}
\DeclareMathOperator{\Supp}{Supp}

\newcommand\independent{\protect\mathpalette{\protect\independenT}{\perp}}
\def\independenT#1#2{\mathrel{\rlap{$#1#2$}\mkern2mu{#1#2}}}
\newcommand{\indep}{\perp\!\!\!\perp}

\usepackage{cleveref}
\crefname{appsec}{appendix}{appendices}
\crefname{appsubsec}{appendix}{appendices}
\crefname{assumption}{assumption}{assumptions}
\crefname{equation}{equation}{equations}
\crefname{exmp}{example}{examples}
\crefname{assN}{assumption}{assumptions}
\crefname{cmt}{comment}{comments}
\crefname{defN}{definition}{definitions}

\usepackage[nolist]{acronym}
\begin{acronym}
  \acro{ATE}{average treatment effect}
  \acro{ATT}{average treatment effect on the treated}
  \acro{CIA}{conditional independence assumption}
  \acro{CI}{confidence interval}
  \acro{CLT}{central limit theorem}
  \acro{DiD}{difference-in-differences}
  \acro{FE}{fixed effects}
  \acro{FWL}{Frisch-Waugh-Lovell}
  \acro{IV}{instrumental variables}
  \acro{LATE}{local average treatment effect}
  \acro{OLS}{ordinary least squares}
  \acro{OVB}{omitted-variable bias}
  \acro{PO}{potential outcomes}
  \acro{RCT}{randomized controlled trial}
  \acro{RD}{regression discontinuity}
  \acro{SUTVA}{stable unit treatment value assignment}
  \acro{TWFE}{two-way fixed effects}
\end{acronym}

\def\inprobHIGH{\,{\buildrel p \over \rightarrow}\,}
\def\inprob{\,{\inprobHIGH}\,}
\def\indistHIGH{\,{\buildrel d \over \rightarrow}\,}
\def\indist{\,{\indistHIGH}\,}

\usepackage[many]{tcolorbox}
\definecolor{main}{HTML}{3F6EA7}
\definecolor{sub}{HTML}{F2F7FD}
\definecolor{sub2}{HTML}{FFF5E8}
\definecolor{mainline}{HTML}{D3DFEE}
\definecolor{warn}{HTML}{C7771B}
\definecolor{warnline}{HTML}{EAD9C0}

\tcbset{
    enhanced,
    breakable,
    colback = white,
    boxrule = 0.5pt,
    arc = 2pt,
    left = 9pt,
    right = 9pt,
    top = 7pt,
    bottom = 7pt,
    before skip = 0.3cm,
    after skip = 0.45cm
}

\newtcolorbox{boxD}{
    colback = sub,
    colframe = mainline,
    borderline west = {2.2pt}{0pt}{main},
    borderline north = {0.6pt}{0pt}{main!25},
    borderline south = {0.6pt}{0pt}{main!15}
}

\newtcolorbox{boxF}{
    colback = sub2,
    colframe = warnline,
    borderline west = {2.2pt}{0pt}{warn},
    borderline north = {0.6pt}{0pt}{warn!25},
    borderline south = {0.6pt}{0pt}{warn!15}
}

\usepackage{colortbl}
\usepackage{tikz}
\usepackage{verbatim}
\usetikzlibrary{positioning}
\usetikzlibrary{snakes}
\usetikzlibrary{calc}
\usetikzlibrary{arrows}
\usetikzlibrary{decorations.markings}
\usetikzlibrary{shapes.misc}
\usetikzlibrary{matrix,shapes,arrows,fit,tikzmark}
